\documentclass[11pt]{article}
\usepackage[a4paper,margin=1in]{geometry}
\usepackage{fontspec}
\usepackage[boldfont=false]{xeCJK}
\setCJKmainfont{FandolSong-Regular.otf}
\usepackage{amsmath}
\usepackage{amssymb}
\usepackage{array}
\usepackage{booktabs}
\usepackage{graphicx}
\usepackage{adjustbox}
\usepackage{enumitem}
\setlist[itemize]{topsep=2pt,partopsep=0pt,parsep=0pt,itemsep=2pt,leftmargin=1.4em}
\setlist[enumerate]{topsep=2pt,partopsep=0pt,parsep=0pt,itemsep=2pt,leftmargin=1.6em}
\usepackage{parskip}
\usepackage{fvextra}
\fvset{breaklines=true,breakanywhere=true,fontsize=\small}
\setlength{\parindent}{0pt}

\begin{document}

\begin{center}
  {\LARGE\bfseries Transport in mammals}\\[0.35em]
  {\large A Level Biology}
\end{center}
\vspace{0.6em}

\section*{The circulatory system}

Mammals have a \textbf{closed double circulation}. "Closed" means the blood stays inside \textbf{blood vessels} 血管 the whole time. "Double" means the blood passes through the \textbf{heart} 心脏 twice for each full trip around the body. This gives two linked loops, so we call it a \textbf{double circulation} 双循环 and the whole thing a \textbf{circulatory system} 循环系统:

\begin{itemize}
\item the \textbf{pulmonary circulation} 肺循环 carries blood from the heart to the lungs and back.

\item the \textbf{systemic circulation} 体循环 carries blood from the heart to the rest of the body and back.

\end{itemize}

Blood travels in this order: \textbf{arteries} 动脉 → \textbf{arterioles} 小动脉 → \textbf{capillaries} 毛细血管 → \textbf{venules} 小静脉 → \textbf{veins} 静脉.

The main vessels are:

\begin{itemize}
\item the \textbf{pulmonary artery} 肺动脉 — carries blood low in \textbf{oxygen} 氧气 from the heart to the lungs.

\item the \textbf{pulmonary vein} 肺静脉 — carries oxygen-rich blood from the lungs back to the heart.

\item the \textbf{aorta} 主动脉 — the big artery that carries oxygen-rich blood from the heart to the body.

\item the \textbf{vena cava} 腔静脉 — the big vein that returns oxygen-poor blood from the body to the heart.

\end{itemize}

\subsection*{How the vessels suit their jobs}

\begin{center}
\adjustbox{max width=\linewidth}{%
\begin{tabular}{lll}
\toprule
\textbf{Vessel} & \textbf{Structure} & \textbf{Function} \\
\midrule
artery & thick wall of \textbf{muscle} 肌肉 and \textbf{elastic} 弹性 fibres; narrow \textbf{lumen} 管腔 (the space inside) & carries blood at high pressure away from the heart; elastic walls stretch and recoil to smooth the flow; muscle controls the flow \\
capillary & wall just one cell thick; very narrow & short distance for exchange of substances between blood and cells \\
vein & thin wall; wide lumen; has \textbf{valves} 瓣膜 & returns blood at low pressure to the heart; valves stop blood flowing backwards \\
\bottomrule
\end{tabular}}
\end{center}

\subsection*{Blood cells}

You should recognise: \textbf{red blood cells} 红细胞 (which carry oxygen), and three white blood cells — \textbf{monocytes} 单核细胞, \textbf{neutrophils} 中性粒细胞 and \textbf{lymphocytes} 淋巴细胞.

\subsection*{Water, plasma and tissue fluid}

Water is the main part of blood. It is a good \textbf{solvent} 溶剂, so it carries dissolved substances, and it has a high \textbf{specific heat capacity} 比热容, so the blood's temperature stays steady.

At the start of a capillary, the high blood pressure pushes liquid (but not the cells or large proteins) out of the \textbf{plasma} 血浆 and through the capillary wall. This liquid around the cells is \textbf{tissue fluid} 组织液. It supplies the cells with oxygen and glucose and carries waste away. Most of it returns to the capillary at the far end, where the pressure is lower.

\section*{Transport of oxygen and carbon dioxide}

\subsection*{Carrying oxygen}

Oxygen is carried by \textbf{haemoglobin} 血红蛋白 in the red blood cells. We show how well haemoglobin holds oxygen with the \textbf{oxygen dissociation curve} 氧解离曲线. This S-shaped graph plots the \textbf{saturation} 饱和度 (how full of oxygen the haemoglobin is) against the \textbf{partial pressure} 分压 of oxygen:

\begin{itemize}
\item where the partial pressure of oxygen is \textbf{high} (in the lungs), haemoglobin loads up and becomes almost fully saturated.

\item where it is \textbf{low} (in respiring tissues), haemoglobin unloads its oxygen for the cells to use.

\end{itemize}

\subsection*{The Bohr shift}

When tissues are very active, they release more \textbf{carbon dioxide} 二氧化碳, which lowers the pH. This makes haemoglobin release oxygen more easily, so the curve moves to the right. This helpful change is the \textbf{Bohr shift} 波尔位移: oxygen is given up exactly where it is most needed.

\subsection*{Carrying carbon dioxide}

A little carbon dioxide dissolves straight into the plasma, but most is carried after a reaction inside the red blood cells:

\begin{enumerate}
\item the \textbf{enzyme} 酶 \textbf{carbonic anhydrase} 碳酸酐酶 speeds up the reaction of carbon dioxide with water to make carbonic acid.

\item the carbonic acid splits into hydrogen ions and \textbf{hydrogencarbonate ions} 碳酸氢根离子.

\item the hydrogencarbonate ions move out into the plasma. This is the main way carbon dioxide is carried.

\item to keep the charge balanced, \textbf{chloride ions} 氯离子 move into the red blood cells. This movement is the \textbf{chloride shift} 氯转移.

\item the hydrogen ions join haemoglobin to form \textbf{haemoglobinic acid} 血红蛋白酸; this mops up the hydrogen ions and keeps the pH steady.

\end{enumerate}

Some carbon dioxide also joins haemoglobin directly to form \textbf{carbaminohaemoglobin} 氨甲酰血红蛋白.

\section*{The heart}

\subsection*{Structure}

The heart has four chambers. The two upper chambers are the \textbf{atria} 心房 (singular: atrium); they have thin walls because they only push blood down into the chambers below. The two lower chambers are the \textbf{ventricles} 心室; they have thick muscular walls because they pump blood out of the heart.

The left ventricle wall is thicker than the right ventricle wall, because the left side must pump blood all the way round the body, while the right side only pumps to the nearby lungs.

\subsection*{The cardiac cycle}

One heartbeat is the \textbf{cardiac cycle} 心动周期. It has two parts: \textbf{systole} 收缩期 (when heart muscle contracts) and \textbf{diastole} 舒张期 (when it relaxes and fills). When a chamber contracts, the pressure inside rises; this pressure change opens and closes the valves so that blood flows one way only.

\subsection*{Controlling the heartbeat}

The heart sets its own rhythm:

\begin{itemize}
\item the \textbf{sinoatrial node} 窦房结 in the right atrium is the pacemaker. It sends out a wave of electrical excitation that spreads across the atria and makes them contract.

\item the \textbf{atrioventricular node} 房室结 picks up the wave, holds it back for a moment (so the atria empty first), then passes it on.

\item the \textbf{Purkyne tissue} 浦肯野组织 carries the wave down and through the ventricle walls, so the ventricles contract from the bottom upwards and push blood out.

\end{itemize}


\end{document}
